\documentclass[12pt,a4paper]{article}

% Paquetes esenciales para idioma y caracteres especiales
\usepackage[utf8]{inputenc}
\usepackage[spanish]{babel}
\usepackage{amsmath}

% Información del documento
\title{Documento de Prueba en \LaTeX}
\author{Usuario de Prueba}
\date{\today}

\begin{document}

\maketitle

\section{Introducción}
Este es un documento de prueba generado para mostrar la estructura básica de un archivo \LaTeX. Es ideal para empezar a familiarizarse con el entorno.

\section{Ecuaciones Matemáticas}
\LaTeX{} es famoso por su excelente capacidad para renderizar fórmulas matemáticas de alta calidad. Aquí tienes un ejemplo clásico:

\begin{equation}
    E = mc^2
\end{equation}

También puedes escribir matemáticas en línea con el texto, como por ejemplo la fracción $\frac{1}{2}$ o la integral $\int_{0}^{1} x \, dx$.

\section{Listas}
Crear listas estructuradas es muy sencillo:

\begin{itemize}
    \item Primer elemento de la lista.
    \item Segundo elemento, que puede ser más largo para ver cómo se ajusta el texto.
    \item Tercer elemento.
\end{itemize}

\end{document}